\subsubsection{Construction via Fourier transform}\label{sss:Four} By the semismallness of $\pi$,  the complex $\calS=\bR\pi_{*}\Ql[\dim\sN]$ is also a perverse sheaf. However, it is not the middle extension from an open subset of $\calN$. There is a notion of Fourier transform for $\Gm$-equivariant sheaves on affine spaces \cite{Laumon}. One can show that $\calS$ is isomorphic to the Fourier transform of $\calS_{\frg}$ and vice versa. The $W$-action on $\calS_{\frg}$ then induces an action of $W$ on $\calS$ by the functoriality of Fourier transform. Taking the stalk of $\calS$ at $e$ we get an action of $W$ on $\cohog{*}{\sB_{e}}$. This action is normalized according to our convention in \S\ref{sss:two actions}.

\subsubsection{Construction via correspondences}  Consider the Steinberg variety $\St_{\frg}=\tg\times_{\frg}\tg$ which classifies triples $(X,B_{1},B_{2})$, where $B_{1},B_{2}$ are Borel subgroups of $G$ and $X\in\Lie B_{1}\cap\Lie B_{2}$. The irreducible components of $\St_{\frg}$ are indexed by elements in the Weyl group: for $w\in W$, letting $\St_{w}$ be the closure of the graph of the $w$-action on $\tg^{\rs}$, then $\St_{w}$ is  an irreducible component of $\St$ and these exhaust all irreducible components of $\St$. The formalism of cohomological correspondences allows us to get an endomorphism of the complex $\calS_{\frg}=\bR\pi_{\frg,*}\Ql[\dim\frg]$ from each $\St_{w}$. It is nontrivial to show that these endomorphisms together form an action of $W$ on $\calS_{\frg}$. The key ingredient in the argument is still the smallness of the map $\pi_{\frg}$. After the $W$-action on $\calS_{\frg}$ is defined, one then define the Springer action on $\cohog{*}{\sB_{e}}$ by either twisting the action of $W$ on the stalk $\calS_{\frg, e}$ by the sign representation as in \S\ref{sss:mid}, or by using Fourier transform as in \S\ref{sss:Four}. We refer to \cite[Remark 3.3.4]{GS} for some discussion of this construction. See also \cite[\S3.4]{CG} for a similar but different construction using limits of $\St_{w}$ in the nilpotent Steinberg variety $\tN\times_{\sN}\tN$. 

Note that the above three constructions all allow one to show that $\End(\calS)\cong\Ql[W]$, hence giving a proof of Theorem \ref{th:BM}.


\subsubsection{Construction via monodromy} We sketch a construction of Slodowy \cite[\S4]{Slo4} which works for $k=\CC$. This construction was conjectured to give the same action of $W$ on $\cohog{*}{\sB_{e}}$ as the one in Theorem \ref{th:Spr corr}.  A similar construction by Rossmann appeared in \cite[\S2]{Rossmann}, in which the author identified his action with that constructed by Kazhdan and Lusztig in \cite{KL80}, and the latter was known to be the same as Springer's action. Thus all these constructions give the same $W$-action as in Theorem \ref{th:Spr corr}.

Let $e\in\calN$ and let $S_{e}\subset \frg$ be a transversal slice to the orbit of $e$. Upon choosing an $\sl_{2}$-triple $(e,h,f)$ containing $e$, there is a canonical choice of such a transversal slice $S_{e}=e+\frg_{f}$, where $\frg_{f}$ is the centralizer of $f$ in $\frg$. Now consider the following diagram where the squares are Cartesian except for the rightmost one
\begin{equation}\label{Slo}
\xymatrix{\sB_{e}\ar[d]\ar@{^{(}->}[r] & \tilS^{\nil}_{e}\ar[d]\ar@{^{(}->}[r] & \tilS_{e}\ar[d]\ar@{^{(}->}[r] & \tg\ar[d]^{\pi_{\frg}}\ar[r]^{\wt\chi} & \frt\ar[d]\\
\{e\}\ar@{^{(}->}[r] & S^{\nil}_{e}\ar@{^{(}->}[r] & S_{e}\ar@{^{(}->}[r] & \frg\ar[r]^{\chi} & \frc}
\end{equation}
Here the rightmost square is \eqref{tgg}, $S^{\nil}_{e}=S_{e}\cap\sN$ and $\tilS_{e}$ and $\tilS^{\nil}_{e}$ are the preimages of $S_{e}$ and $S^{\nil}_{e}$ under $\pi_{\frg}$. Let $V_{e}\subset S_{e}$ be a small ball around $e$ and let  $V_{0}\subset \frc$ be an even smaller ball around $0\in\frc$. Let $U_{e}=V_{e}\cap \chi^{-1}(V_{0})\subset S_{e}$. Then the diagram \eqref{Slo} restricts to a diagram
\begin{equation}\label{Slo U}
\xymatrix{\sB_{e}\ar[d]\ar@{^{(}->}[r] & \tilU^{\nil}_{e}\ar[d]\ar@{^{(}->}[r] & \tilU_{e}\ar[d]\ar[r]^{\wt\chi_{e}} & \tilV_{0}\ar[d]\\
\{e\}\ar@{^{(}->}[r] & U^{\nil}_{e}\ar@{^{(}->}[r] & U_{e}\ar[r] & V_{0}}
\end{equation} 
Here $U^{\nil}_{e}=U_{e}\cap\sN$, and $\tilV_{0}, \tilU_{e}$ and $\tilU^{\nil}_{e}$ are the preimages of $V_{0}, U_{e}$ and $U^{\nil}_{e}$ under the vertical maps in \eqref{Slo}. The key topological facts here are
\begin{itemize}
\item The inclusion $\sB_{e}\incl \tilU^{\nil}_{e}$ admits a deformation retract, hence it is a homotopy equivalence;
\item The map $\wt\chi_{e}: \tilU_{e}\to \tilV_{0}$ is a trivializable fiber bundle (in the sense of differential topology).
\end{itemize}
Now a general fiber of $\wt\chi_{e}$ admits a homotopy action of $W$ by the second point above because the rightmost  square in \eqref{Slo U} is Cartesian over $V_{0}\cap\frc^{\rs}$ and the map $\tilV_{0}\to V_{0}$ is a $W$-torsor over $V_{0}\cap\frc^{\rs}$. By the first point above, $\sB_{e}$ has the same homotopy type with $\tilU^{\nil}_{e}=\wt\chi^{-1}_{e}(0)$, hence $\sB_{e}$ also has the same homotopy type as a general fiber of $\wt\chi_{e}$ because $\wt\chi_{e}$ is a fiber bundle. Combining these facts, we get an action of $W$ on the {\em homotopy type of $\sB_{e}$}, which is a stronger structure than an action of $W$ on the cohomology of $\sB_{e}$. A consequence of this construction is that the $W$-action on $\cohog{*}{\sB_{e}}$ in Theorem  \ref{th:Spr corr} preserves the ring structure.


 

\subsubsection{Proof of Theorem \ref{th:Spr corr} assuming Theorem \ref{th:BM}}
We decompose the perverse sheaf $\calS$ into isotypical components under the $W$-action
\begin{equation*}
\calS=\bigoplus_{\chi\in\Irr(W)}V_{\chi}\otimes\calS_{\chi}
\end{equation*}
where $V_{\chi}$ is the space on which $W$ acts via the irreducible representation $\chi$, and $\calS_{\chi}=\Hom_{W}(V_{\chi},\calS)$ is a perverse sheaf on $\sN$. Since $\End(\calS)\cong\Ql[W]$, we conclude that
each $\calS_{\chi}$ is nonzero and that
\begin{equation}\label{HomSS}
\Hom(\calS_{\chi},\calS_{\chi'})=\begin{cases}\Ql & \chi=\chi' \\ 0 & \textup{otherwise}\end{cases}
\end{equation}
The decomposition theorem \cite[Th 6.2.5]{BBD} implies that each $\calS_{\chi}$ is a semisimple perverse sheaf. Therefore \eqref{HomSS} implies that $\calS_{\chi}$ is simple. Hence $\calS_{\chi}$ is of the form $\IC(\ov\calO, \calL)$ where $\calO\subset\sN$ is a nilpotent orbit and $\calL$ is an irreducible  $G$-equivariant local system on $\calO$. Moreover, since $\Hom(\calS_{\chi},\calS_{\chi'})=0$ for $\chi\neq\chi'$, the simple perverse sheaves $\{\calS_{\chi}\}_{\chi\in\Irr(W)}$ are non-isomorphic to each other. This proves part (3) of Theorem \ref{th:Spr corr} by interpreting the right side of \eqref{Spr corr} as the set of isomorphism classes of irreducible $G$-equivariant local systems on nilpotent orbits. If $\calS_{\chi}=\IC(\ov\calO,\calL)$ and $e\in\calO$, the semismallness of $\pi$ allows us to identify the stalk $\calL_{e}$ with an $A_{e}$-isotypic subspace of $\cohog{2d_{e}}{\sB_{e}}$. This proves part (2) of Theorem \ref{th:Spr corr}. \qed



We give some further examples of the Springer correspondence.



\subsubsection{Type $A$}\label{sss:A Spr corr} When $G=\SL_{n}$, all $A_{e}$ are trivial. The Springer correspondence sets a bijection between irreducible representations of $W=S_{n}$ and nilpotent orbits of $\frg=\sl_{n}$, both parametrized by partitions of $n$. In \S\ref{sss:Spr type A} we have seen that if $e$ has Jordan type $\l$, the top dimensional cohomology $\cohog{2d_{e}}{\sB_{e}}$ has a basis indexed by the standard Young tableaux of $\l$, the latter also indexing a basis of the irreducible representation of $S_{n}$ corresponding to the partition $\l$. For example, for $G=\SL_{3}$, the Springer correspondences reads
\begin{itemize}
\item trivial representation $\bij$ regular orbit, partition $3=3$;  
\item two-dimensional representation $\bij$ subregular orbit, partition $3=2+1$; 
\item sign representation $\bij$ \{0\}, partition $3=1+1+1$. 
\end{itemize}


\subsubsection{The subregular orbit and the reflection representation}\label{sss:subreg Spr corr} Consider the case $e$ is a subregular nilpotent element. In this case, the component group $A_{e}$ can be identified with the automorphism group of the Dynkin diagram of $G'$ introduced in \S\ref{sss:subreg} (see \cite[\S7.5, Proposition]{Slodowy}). After identifying the irreducible components of $\sB_{e}$ with the vertices of the Dynkin diagram of $G'$, the action of $A_{e}$ on $\cohog{2}{\sB_{e}}$ is by permuting the basis given by irreducible components in the same way as its action on the Dynkin diagram of $G'$. For example, when $G=G_{2}$,  we may write $\sB_{e}=C_{1}\cup C_{2}\cup C_{3}\cup C_{4}$ with $C_{1}$, $C_{2}$, $C_{3}$ each intersecting $C_{4}$ in a point and otherwise disjoint. The group $A_{e}$ is isomorphic to $S_{3}$, and its action on $\cohog{2}{\sB_{e}}$ fixes the fundamental class of $C_{4}$ and permutes the fundamental classes of $C_{1}, C_{2}$ and $C_{3}$.

Note that $\cohog{2}{\sB_{e}}^{A_{e}}$ always has dimension $r$, the rank of $G$. In fact, as a $W$-module, $\cohog{2}{\sB_{e}}^{A_{e}}$ is isomorphic to the reflection representation of $W$ on $\frt^{*}$. In other words, under the Springer correspondence, the pair $(e=\textup{subregular}, \rho=1)$ corresponds to the reflection representation of $W$.





